\documentclass[11pt]{article}
\usepackage[margin=0.78in]{geometry}
\usepackage{booktabs}
\usepackage{graphicx}
\usepackage[section]{placeins}
\usepackage{hyperref}
\usepackage{xcolor}
\usepackage{enumitem}
\usepackage{longtable}
\usepackage{tabularx}
\usepackage{microtype}
\usepackage{fontspec}
\setmainfont{TeX Gyre Heros}
\setsansfont{TeX Gyre Heros}
\definecolor{linkblue}{HTML}{0068A8}
\hypersetup{colorlinks=true,linkcolor=linkblue,urlcolor=linkblue}
\setlist{nosep,leftmargin=1.4em}
\setlength{\parskip}{0.48em}
\setlength{\parindent}{0pt}
\setlength{\emergencystretch}{2em}
\newcommand{\figdir}{artifacts/reports/native-sft-study-20260914/figures}
\newcommand{\code}[1]{\texttt{#1}}
\title{Teaching Qwen to Edit and Verify TRExFitter Configurations\\
\large Native-tool supervised fine-tuning study}
\author{ATLAS H$\rightarrow\gamma\gamma$ coding-agent project}
\date{14 September 2026}

\begin{document}
\maketitle

\begin{abstract}
This report studies whether supervised fine-tuning (SFT) can teach small and
mid-size Qwen models to behave as reliable coding agents for TRExFitter
configuration work.  The target behavior is a short, auditable workflow:
understand an analyst's requested change, modify \code{analysis.config} with
ordinary coding tools, run an independent Python verifier, and report only an
observed result.  The public dataset contains 167 structured trajectories over
156 logical tasks.  Every deterministic demonstration was replayed against a
hidden oracle before admission.  Untouched public checkpoints and their SFT
counterparts are evaluated under identical prompts, tools, decoding, turn
limits, and held-out tasks.  The results show that SFT can teach parts of the
tool protocol, but low validation loss alone does not imply reliable end-to-end
agent behavior; the best hosted coding-agent control remains substantially
ahead on this first 36-task benchmark.
\end{abstract}

\section{The problem in plain language}

An ATLAS analysis configuration is text, but editing it safely is not merely a
text-generation problem.  A useful agent must choose a scientifically
meaningful setting, place it in the correct named block, preserve unrelated
settings and references, and prove that the resulting file remains valid.
Simple exact-string comparison is too brittle: harmless formatting can differ,
and some expressions can be semantically equivalent.  Conversely, a file that
parses can still implement the wrong analysis decision.

This study deliberately starts with a conservative SFT target.  A trajectory
passes only when the requested state is present, the file is statically valid,
unrelated content is preserved, and a successful verifier invocation occurs
after the final edit.  Running TRExFitter and comparing produced histograms is a
stronger semantic-equivalence level reserved for reinforcement learning (RL).

\section{Dataset: analysis decisions, not arbitrary prose}

The public release is
\href{https://huggingface.co/datasets/cxyang-ucb/hyy-sft}{\code{cxyang-ucb/hyy-sft}}.
Its immutable Hugging Face revision is:
\begin{center}
\small\texttt{ea52627d07328d0a554886989dda96e444921ee6}
\end{center}
It has 129 training and 38 validation trajectories.  These correspond to 120
training and 36 validation \emph{logical tasks}: nine tasks have both a
deterministic trajectory and a real teacher rollout in training, and two have
both forms in validation.  The mixture is:

\begin{itemize}
  \item 156 deterministic generic-tool demonstrations, all independently replayed;
  \item 11 real coding-agent trajectories from Codex and OpenCode, retained with provenance;
  \item nine kinds of analysis work, from single settings through whole-block reconstruction.
\end{itemize}

\begin{figure}[ht]
\centering
\includegraphics[width=\textwidth]{\figdir/dataset-composition.png}
\caption{Logical-task coverage.  ``Complexity'' is a construction label: 1 is a
single local decision; 4 can require coordinated settings or block reconstruction.}
\end{figure}

Settings with a closed choice set---for example a Boolean, an enumerated fit
mode, or a valid region reference---are generated from documented policies.
Open-ended fields such as titles are varied as natural language.  Tasks include
additions, repairs, coordinated changes, no-op recognition, reference repair,
and reconstruction of a removed Fit, Region, Sample, or NormFactor block from a
natural-language specification.  This is intended to teach both syntax and the
human analysis decisions encoded by the file.

\subsection{What one training trajectory teaches}

A representative demonstration asks the agent to add \code{FitBlind: TRUE}.
The assistant first calls the generic \code{apply\_patch} tool with a minimal
diff.  After receiving the tool result, it calls \code{shell} with
\code{python -m trex\_fitter.config\_verify analysis.config --actions n
--evidence-level S}.  Only after observing exit code zero does it report that
validation passed.  A held-out example follows the same protocol but asks for
\code{UseMinos: mu\_H}; its answer is never used in training.

\begin{verbatim}
assistant -> apply_patch:
  Fit: "fit"
+   FitBlind: TRUE
    FitType: SPLUSB

tool -> assistant: Patch applied successfully.
assistant -> shell: python -m trex_fitter.config_verify ...
tool -> assistant: Exit code: 0; VALID; COMPATIBLE
assistant: requested state is present and validation passed
\end{verbatim}

These are ordinary harness capabilities---read/search a file, patch it, and run
a bounded shell command.  They are neither MCP calls nor a TRExFitter-specific
model-facing API.  The Python verifier is domain-specific, but the model invokes
it through the same shell tool a coding agent already has.

\section{Why the trajectories are valid}

Each deterministic row is admitted only after an executable replay starts from
its recorded fixture, performs the recorded patch, runs the exact recorded
verifier command, checks the stored tool observations, compares the final file
byte-for-byte with a hidden golden file, checks the semantic change contract,
and audits preservation plus the final response.  All 156 deterministic rows
passed this gate.  Teacher rollouts use a separate scorer and must contain a
successful observed verifier call after their final edit.

This makes the golden response executable evidence, not merely plausible model
text.  It does not prove that an analyst would endorse every scientific choice;
human review remains required before production use.

\section{Model-neutral records and Qwen-native rendering}

Canonical rows store structured \code{messages} and optional OpenAI-style
\code{tools}.  They do not store Qwen control tokens, rendered strings, token
IDs, labels, or hidden reasoning.  VERL passes the structures to the exact
checkpoint tokenizer.  That tokenizer's chat template therefore owns the
model-specific function-call syntax.

Before training, every row passes a tokenizer/template gate.  For Qwen3.5 the
largest training trajectory is 10,335 tokens and the largest validation
trajectory is 8,158 tokens; the study uses a hard 12,288-token limit and rejects
truncation.  The gate also proves that VERL's turn-wise rendering equals the
tokenizer's full-conversation rendering, that only assistant spans receive
loss, and that thinking is disabled consistently.  Qwen3.5 is capable of a
thinking mode, but this study does not supervise hidden chain-of-thought: the
template is called with \code{enable\_thinking=False}, empty think delimiters
are masked, and only observable tool behavior and final answers are learned.

\section{Training design}

Training uses VERL's documented
\href{https://verl.readthedocs.io/en/latest/examples/config.html}{SFT trainer}
(\code{verl.trainer.sft\_trainer}), BF16, FSDP, gradient
checkpointing, one epoch, seed 1, and assistant-only loss.  ``Full SFT'' updates
all model parameters.  ``LoRA'' (low-rank adaptation) trains small low-rank
matrices while leaving the original weights frozen, reducing memory and
checkpoint size.  The main ablation changes method, learning rate, or LoRA
rank while holding data and evaluation fixed.

\begin{figure}[ht]
\centering
\includegraphics[width=0.96\textwidth]{\figdir/training-loss.png}
\caption{Loss on assistant tokens during one epoch.  It measures next-token
prediction on demonstrations, not task success.  Short-run oscillation is
expected because batches have different task and sequence mixtures.}
\end{figure}

\begin{figure}[ht]
\centering
\includegraphics[width=0.9\textwidth]{\figdir/validation-loss-ablation.png}
\caption{End-of-epoch validation loss.  Lower is better for reproducing held-out
assistant trajectories; verifier-scored execution remains the decision metric.}
\end{figure}

\section{Evaluation: the controlled base comparison}

``Base'' means the untouched public instruct checkpoint with no project SFT.
It does not mean a raw, instruction-free pretraining checkpoint.  For every
\emph{Qwen base/SFT pair}, the evaluator uses the same 36 held-out logical
tasks, structured prompt, four generic tools, four-turn limit, 768-token
generation limit, temperature zero, seed 1, isolated temporary workspace, and
independent scorer.  Thus the only intended difference inside a Qwen pair is
project SFT.  Hosted controls use their harness-native prompt and filesystem
interaction, but share the logical tasks, initial files, allowed operations,
verifier, and scorer; they are informative references rather than a causal
model-only comparison.

A \textbf{strict pass} requires all of the following: requested state present;
static config validity; preservation of unrelated state; any required evidence;
and a successful verifier call after the last patch.  The denominator is always
36.  Error bars in the leaderboard are 95\% Wilson binomial intervals; with
only 36 tasks they are necessarily wide.

\begin{figure}[ht]
\centering
\includegraphics[width=0.94\textwidth]{\figdir/base-vs-sft.png}
\caption{The primary causal comparison.  Each line connects an untouched
checkpoint to its trained counterpart under matched evaluation settings.}
\end{figure}

\begin{table}[ht]
\centering
\small
\begin{tabular}{llrrr}
\toprule
Checkpoint & Project training & Base & After SFT & Change \\
\midrule
Qwen2.5-Coder-0.5B & Full & 0/36 & 0/36 & 0.0 pp \\
Qwen3.5-0.8B & Full & 0/36 & 4/36 & +11.1 pp \\
Qwen3.5-0.8B & LoRA & 0/36 & 1/36 & +2.8 pp \\
Qwen3.5-2B & Full & 0/36 & 2/36 & +5.6 pp \\
Qwen3.5-4B & LoRA & 0/36 & 9/36 & +25.0 pp \\
Qwen3.5-9B & LoRA & 0/36 & 13/36 & +36.1 pp \\
Qwen3.5-35B-A3B & LoRA & 0/36 & 0/36 & 0.0 pp \\
\bottomrule
\end{tabular}
\caption{Exact matched results.  ``pp'' means percentage points.  Base is the
same public instruct checkpoint before any project SFT.}
\end{table}

\begin{figure}[ht]
\centering
\includegraphics[width=0.95\textwidth]{\figdir/strict-task-success.png}
\caption{Verifier-backed task success.  Hosted controls are references, not a
claim of a universal state-of-the-art ranking; their model and default effort
are printed explicitly.}
\end{figure}

\section{Results and observations}

\textbf{Small-model result.}  Qwen2.5-Coder-0.5B scored 0/36 both before and
after this one-epoch SFT.  Its base model often printed JSON describing a
nonexistent ``validate'' function instead of calling the supplied coding tools;
the trained checkpoint still did not cross the strict end-to-end gate.  This is
useful negative evidence: a falling training loss is not enough at this scale
and data volume.

\textbf{0.8B controlled result.}  Untouched Qwen3.5-0.8B scored 0/36.  Full SFT
reached 4/36 (11.1\%) and invoked the verifier on 16/36 tasks; LoRA reached
1/36 (2.8\%).  The base occasionally produced a correct final file but never
observed a successful verifier, so it failed the agent protocol.  Full SFT's
gain is encouraging but uncertain: for four discordant wins over base, a
two-sided exact McNemar test gives $p=0.125$.  More held-out tasks and seeds are
needed before claiming a stable effect size.

\textbf{Training method.}  On 0.8B, full SFT obtained both lower validation loss
(0.0452) and higher strict success than the selected LoRA run (loss 0.0627;
1/36).  Two additional LoRA settings were worse in validation loss: rank 16 at
learning rate $5\times10^{-5}$ gave 0.1200, and rank 8 at $10^{-4}$ gave
0.0999.  This is a small one-seed ablation, not a general verdict against LoRA.

\textbf{Scale.}  Qwen3.5-2B full SFT reached validation loss 0.0374 and improved
strict success from 0/36 to 2/36 (5.6\%).  Qwen3.5-4B LoRA reached loss 0.0433
and improved from 0/36 to 9/36 (25.0\%).  Qwen3.5-9B LoRA reached loss 0.0427
and improved from 0/36 to 13/36 (36.1\%), the strongest open-weight result in
this study.  Loss ordering therefore should not be substituted for verifier
accuracy.  The untouched 4B and 9B checkpoints each produced ten correct final
configs yet scored 0/36 strictly because neither completed a successful
post-edit verifier call.  At these sizes, protocol adherence---not only text
editing---is therefore the dominant measured failure.
Because the 9B base had no wins, its 13 discordant SFT wins give a two-sided
exact McNemar $p=0.00024$ (unadjusted for the several model comparisons).
This supports a paired effect on this benchmark, while the single training
seed still limits generalization.

\textbf{Largest-model result.}  Qwen3.5-35B-A3B LoRA had validation loss
0.2105 and scored 0/36 both before and after SFT.  Correct final configurations
decreased from 13/36 to 10/36, and neither side observed a successful verifier.
Thus this run provides no evidence of an SFT benefit.  It is not evidence that
35B models are intrinsically worse: the mixture-of-experts architecture,
single LoRA setting, one epoch, and small dataset prevent a controlled
parameter-scaling inference.

\textbf{Hosted controls.}  Under their registry defaults, GPT-6-Astra at low
effort scored 22/36 (61.1\%) and GPT-5.6-Sol at medium effort scored 31/36
(86.1\%).  Effort differs and is a real confound, so the two rows should not be
read as a pure model-generation comparison.  They demonstrate the remaining
gap to mature hosted coding agents under the same task and verifier contract.

\begin{figure}[ht]
\centering
\includegraphics[width=0.96\textwidth]{\figdir/outcome-gates.png}
\caption{Diagnostic gates distinguish editing failure from protocol failure.
``Verifier observed'' means a successful verifier run occurred after the final
edit; merely ending with a valid file is insufficient.}
\end{figure}

\begin{figure}[ht]
\centering
\includegraphics[width=\textwidth]{\figdir/category-success.png}
\caption{Category breakdown.  Each category contains four held-out tasks, so
one task changes a cell by 25 percentage points; interpret this as diagnosis,
not a precise ranking.}
\end{figure}

\FloatBarrier
\subsection{Exploratory temperature study}

Temperature controls sampling randomness: zero selects greedily, while larger
values flatten the next-token distribution and make alternate actions more
likely.  The best 0.8B full-SFT checkpoint was evaluated at 0, 0.2, and 0.7;
nonzero runs use top-$p=0.95$ and three seeds.  Temperature zero is deterministic
under this runtime, so it is run once.  These decoding runs do not replace the
temperature-zero base/SFT comparison.

Greedy decoding passed 4/36 tasks.  At temperature 0.2 the three seeds passed
4, 4, and 5 tasks (mean 4.33); at 0.7 they passed 6, 3, and 1 (mean 3.33).
The initially promising 6/36 result at 0.7 was therefore not stable.  Mean
verifier observations were 14/36 at 0.2 and 15/36 at 0.7, versus 16/36 under
greedy decoding.  On this small sweep, extra randomness adds variance without
a reproducible accuracy benefit.

\begin{figure}[ht]
\centering
\includegraphics[width=0.82\textwidth]{\figdir/temperature-study.png}
\caption{Strict success and verifier use as decoding becomes more random.}
\end{figure}

\section{A100 throughput study}

The strong-scaling experiment holds Qwen3.5-0.8B, full SFT, data order, global
batch 16, maximum length, optimizer, and six measured steps fixed while changing
only the number of A100 GPUs.  The first compile/warm-up step is excluded.
Throughput means total non-padding training-sequence tokens divided by wall-clock
time; it is not generated inference tokens per second or assistant-only tokens.
Scaling efficiency is measured throughput divided by ideal linear scaling from
the one-GPU point.

Measured throughput was 2,846, 3,315, 3,833, and 1,684 sequence tokens/s on
1, 2, 4, and 8 A100s, respectively.  The 4-GPU point is only 34\% of ideal
linear scaling; the 8-GPU point is slower than one GPU.  These are short,
single-run measurements, but they make the practical conclusion clear: adding
GPUs to this fixed small-model workload is counterproductive without increasing
useful work per synchronized step or improving sequence-length balance.

\begin{figure}[ht]
\centering
\includegraphics[width=0.9\textwidth]{\figdir/a100-throughput-scaling.png}
\caption{Measured strong scaling.  Small models and a tiny fixed global batch
can become communication- or input-limited, so more GPUs need not be economical.}
\end{figure}

\section{What the study establishes---and what it does not}

The strongest positive result is behavioral: SFT moved several Qwen3.5 models
from never completing a verifier-backed task to sometimes performing the full
patch--verify--report protocol, with the 9B LoRA run improving from 0/36 to
13/36.  The strongest caution is equally important: assistant-token loss can
look excellent while strict task accuracy remains low, and the 35B run did not
improve at all.  The benchmark therefore must remain executable and matched to
an untouched base checkpoint.

Limitations include one training seed, only 36 evaluated logical tasks, one
epoch, and different architecture families at 0.5B.  The primary base/SFT
comparison uses greedy decoding; only the 0.8B SFT temperature sweep has three
stochastic seeds.  The
35B-A3B model is a mixture-of-experts model
(\href{https://huggingface.co/Qwen/Qwen3.5-35B-A3B}{35B total and 3B activated
parameters according to its model card}), so its result is not a clean
dense-model scaling point.  Hosted controls use different underlying models and default reasoning
efforts.  No result yet tests full TRExFitter execution or histogram-level
semantic equivalence.  Older artifact studies used a different dataset/tool
contract and are archived rather than merged into this leaderboard.

\section{Recommended next experiments}

\begin{enumerate}
  \item Expand the held-out benchmark and run at least three training and
  decoding seeds; report paired confidence intervals.
  \item Increase coverage of failed protocol states, especially verifier use,
  malformed tool arguments, and recovery after a tool error.
  \item Evaluate intermediate checkpoints so task accuracy, not only loss, can
  select duration and learning rate.
  \item Add a controlled temperature sweep to the best SFT checkpoint while
  keeping the untouched-base comparison at identical temperatures.
  \item Move to VERL RL only after the SFT baseline is stable.  Reward static
  deterministic correctness first, then add sandboxed TRExFitter execution and
  histogram equivalence as a higher evidence level.
\end{enumerate}

\section*{Definitions}

\begin{description}
\item[SFT (supervised fine-tuning)] Training a model to predict demonstrated assistant actions and answers.
\item[LoRA] A memory-efficient SFT method that learns low-rank weight updates while freezing the base weights.
\item[Assistant-only loss] Prediction error is computed on assistant tokens; system, user, and tool-result tokens provide context but are not targets.
\item[Tool call] Structured model output asking the harness to read, search, edit, or execute; the harness performs the operation and returns an observation.
\item[Verifier] Independent Python code that parses the config, checks the requested contract and preservation, and returns explicit evidence.
\item[Strict pass] Correct requested edit, valid and preserved config, plus a successful verifier invocation after the last patch.
\item[Validation loss] Token-prediction error on trajectories excluded from weight updates; it is not executed-task accuracy.
\item[Hosted control] A reference run from a hosted coding agent evaluated on the same task contract; it is not automatically a universal SOTA claim.
\item[Tokens/s] Here, non-padding SFT sequence tokens processed per wall-clock second across all GPUs.
\end{description}

\section*{Reproducibility pointers}

Dataset: \href{https://huggingface.co/datasets/cxyang-ucb/hyy-sft}{cxyang-ucb/hyy-sft}.\\
Code branch: \code{hyy-native-verl-sft-study}.\\
Plot data: \code{artifacts/reports/native-sft-study-20260914/*.csv}.\\
Per-task results: \code{artifacts/evaluation/*/results.jsonl}.\\
Training metrics and checkpoints: \code{artifacts/checkpoints/verl-*}.

\end{document}
